% Section 10.4, from lectures/L10/appendix/c_exercises.md.

\section{Exercises}
\label{c10:sec:exercises}\label{c10:app:c}

\begin{aside}
Nine, ending with the Cross-check. Do them in order; Exercise~\ref{c10:ex:6} onwards need the target
running.

Where an exercise can be checked mechanically, a \textbf{Check yourself} line says how.

\textbf{Solutions.} The exercises that are not Code are answered in \secref{sol:sec:c10}. The Code
exercises have no solution: each is checked by running it, as its Check yourself line says.
\end{aside}

% ------------------------------------------------------------------------------------------
\exercise[fillaccent2]{Bus, device, driver}{Recall}
\label{c10:ex:1}

\xpart{a} Define each of the three, in one sentence, without naming a physical bus.

\xpart{b} The platform bus corresponds to nothing on any schematic. What is it for, and what
property do its devices share?

\xpart{c} On the target, \file{/sys/bus/platform/devices/c000000.qa-dev} exists before any driver is
loaded. Who created it, from what, and when?

\xpart{d} Does a device have to exist before its driver is registered? Describe both orders and say
what the driver author has to do differently for each.

\xpart{e} A device can appear under both \file{/sys/bus} and \file{/sys/class}. What is the
difference in what they group by, and which one would a program looking for ``all the serial
ports'' use?

% ------------------------------------------------------------------------------------------
\exercise[fillaccent2]{Probe and remove}{Recall}
\label{c10:ex:2}

\xpart{a} \code{probe} can be called more than once for one loaded module. When, and what does that
rule out about where a driver keeps its state?

\xpart{b} A driver keeps \code{static void __iomem *regs;} at file scope and sets it in
\code{probe}. It works. Describe the machine on which it stops working.

\xpart{c} \code{remove} returns \code{void} in this kernel. Why is it not allowed to fail?

\xpart{d} What is \code{-EPROBE_DEFER} for, what does the core do with it, and what is the
alternative that makes your driver depend on link order?

\xpart{e} In \code{probe}, what should be done last, and why? Relate it to \S\ref{c5:app:a3}'s rule
about \code{cdev_add}.

% ------------------------------------------------------------------------------------------
\exercise[fillaccent3]{Reading a device tree}{Hand calculation}
\label{c10:ex:3}

\begin{console}
/ {
        #address-cells = <2>;
        #size-cells = <2>;

        bus@10000000 {
                #address-cells = <1>;
                #size-cells = <1>;
                ranges = <0x1000 0x0 0x10000000 0x100000>;

                thing@2000 {
                        compatible = "acme,thing-v2", "acme,thing";
                        reg = <0x2000 0x40>;
                        interrupts = <0 45 1>;
                };
        };
};
\end{console}

\xpart{a} How many cells does \code{thing}'s \code{reg} occupy, and how do you know?

\xpart{b} How many cells does the bus's \code{ranges} occupy? Account for each.

\xpart{c} What physical address do \code{thing}'s registers start at? Show the arithmetic, and note
that this \code{ranges} does not start its child range at zero.

\xpart{d} \code{compatible} has two entries. Which does the kernel prefer? What is the second for?

\xpart{e} Decode \code{interrupts = <0 45 1>} for a GIC. What hardware interrupt number is that, and
what trigger type?

\xpart{f} A driver's \code{of_match_table} contains only \code{"acme,thing"}. Does it bind? What if
it contained only \code{"acme,thing-v3"}?

% ------------------------------------------------------------------------------------------
\exercise[fillaccent3]{What the helpers do}{Hand calculation}
\label{c10:ex:4}

For each, say what it reads out of the device tree and what it returns on failure.

\xpart{a} \code{devm_platform_ioremap_resource(pdev, 0)}

\xpart{b} \code{platform_get_irq(pdev, 0)}

\xpart{c} \code{of_property_read_u32(node, "fifo-depth", &v)}

Then:

\xpart{d} One of the three returns a negative error code as an \code{int}, one returns an error
packed into a pointer, and one returns zero for success. Say which is which and how each is checked.

\xpart{e} A driver checks \code{if (!irq)} after \code{platform_get_irq}. On which boards does that
work, and on which does it silently do the wrong thing?

\xpart{f} \code{devm_platform_ioremap_resource} replaces four separate calls from \S\ref{c6:app:b7}.
Name all four.

% ------------------------------------------------------------------------------------------
\exercise[fillaccent4]{What belongs in a device tree}{Design}
\label{c10:ex:5}

For each, say whether it belongs in the device tree, in a module parameter, or neither, and why.

\xpart{a} The physical address of the device's registers.

\xpart{b} Which interrupt it raises.

\xpart{c} How deep the driver's software FIFO should be.

\xpart{d} The number of hardware channels the chip has.

\xpart{e} Whether to enable verbose debugging.

\xpart{f} The crystal frequency feeding the device.

\xpart{g} The name of the \file{/dev} node to create.

Then: \textbf{h)} state the rule in one sentence. \textbf{i)} A device tree is an ABI. What follows
for anything you put in one, and why is that a stronger constraint than it is for a module
parameter?

% ------------------------------------------------------------------------------------------
\exercise{A driver that is told nothing}{Code}
\label{c10:ex:6}

Rewrite your Chapter~\ref{c8:ch} driver as \file{lectures/L10/lab/qa_platform.c}, a platform driver.
\begin{itemize}
  \item An \code{of_device_id} table matching \code{QA_DEV_DT_COMPATIBLE}, and
    \code{MODULE_DEVICE_TABLE}.
  \item \code{probe} uses \code{devm_kzalloc} for a per-device private structure, and
    \code{platform_set_drvdata} to attach it.
  \item Registers from \code{devm_platform_ioremap_resource}, interrupt from
    \code{platform_get_irq}, handler from \code{devm_request_irq} named with
    \code{dev_name(&pdev->dev)}.
  \item Check the identity register and fail with \code{-ENODEV} if it is wrong.
  \item Two read-only sysfs attributes, \code{interrupts} and \code{samples}, exposed through the
    driver's \code{.dev_groups} so that the core adds and removes them for you.
  \item \code{remove} stops the device. It must do nothing else.
  \item Use \code{dev_info} and \code{dev_err}, not \code{pr_info}.
\end{itemize}
\textbf{There must be no address and no interrupt number anywhere in the file}, and no \code{goto}
ladder.

\checkyourself \code{make test L=L10} reports \textbf{PASSED}, with nineteen checks. \code{grep -c
'0xC000'} on your source returns 0.

\xpart{a} Compare the line count against your Chapter~\ref{c8:ch} driver. Where did the difference
go?

\xpart{b} Your Chapter~\ref{c6:ch} driver named its own region and the entry in \file{/proc/iomem}
read \code{qa_mmio}. What does it read now, and who chose that name?

\xpart{c} Change one character of the compatible string in your \code{of_match_table} and reload.
What happens, what does \code{dmesg} say, and how would you have diagnosed it without knowing what
you had changed?

% ------------------------------------------------------------------------------------------
\exercise{Bind, unbind, and what \code{devm_} is attached to}{Code}
\label{c10:ex:7}

With your driver loaded and bound:

\xpart{a} Unbind it by hand:

\begin{shell}
echo c000000.qa-dev > /sys/bus/platform/drivers/qa_platform/unbind
\end{shell}

\noindent Check \code{lsmod}, \file{/proc/iomem}, \file{/proc/interrupts} and the device's sysfs
attributes. Which of them changed, and is the module still loaded?

\xpart{b} From that result, state precisely what event releases \code{devm_} resources. Is it module
unload? Give the evidence.

\xpart{c} Bind it again. Did \code{probe} run? How do you know from the sysfs attributes alone?

\xpart{d} Now unload the module without unbinding first. What happens, and in what order?

\xpart{e} Why is \code{unbind} useful in real work? Give two uses that are not demonstrations.

\xpart{f} A process has your device open when you unbind. What keeps the memory alive, and what is
now dangerous about any pointer your file operations still hold?

% ------------------------------------------------------------------------------------------
\exercise[fillaccent4]{Autoloading}{Design}
\label{c10:ex:8}

This target has no udev, so modules are loaded by hand. The rest of the mechanism is present.

\xpart{a} Name the four steps between a device appearing and its module being loaded, and say which
one is missing here.

\xpart{b} Which of the four does \code{MODULE_DEVICE_TABLE} participate in? What breaks if you omit
it, and does the driver still work when loaded by hand?

\xpart{c} Find the \code{MODALIAS} or \code{OF_COMPATIBLE_0} line in your device's \code{uevent}.
Where did that string come from?

\xpart{d} Run \code{modinfo} on your \code{.ko} and find the alias it advertises. Which file would
\code{depmod} have collected it into?

\xpart{e} A USB device you plug in is autoloaded by the same four steps. Where does the string its
\code{uevent} carries come from, and why does a device soldered to the board need a device tree to
supply it?

% ------------------------------------------------------------------------------------------
\exercise[fillaccent]{An address derived three ways}{Cross-check}
\label{c10:ex:9}

You have now computed this address by hand once, in Chapter~\ref{c6:ch}. This time you compute it,
the harness computes it, and the kernel computes it, and all three must agree.

\xpart{a} \textbf{By hand, from the raw bytes.} On the target:

\begin{shell}
dt=/sys/firmware/devicetree/base
node=$(find $dt -name 'qa-dev@*' -type d | head -1)
od -An -tx1 -v "$node/reg"
od -An -tx1 -v "$(dirname "$node")/ranges"
od -An -tx1 -v "$node/interrupts"
\end{shell}

\noindent Write out every cell, four bytes each, most significant first. Then compute the physical
address, the size, and the GIC hardware interrupt number.

\xpart{b} \textbf{From the kernel's own naming.} Look at \code{ls /sys/bus/platform/devices/ | grep
qa-dev}. The device's name contains a number. Which of your computed values is it, and what does
that tell you about when the translation happened?

\xpart{c} \textbf{From the driver.} Have your \code{probe} print what it was handed.
\code{platform_get_resource(pdev, IORESOURCE_MEM, 0)} gives a \code{struct resource} with
\code{start} and \code{end}; print both, and print the value \code{platform_get_irq} returned.

\xpart{d} \textbf{Reconcile.} All three should agree on the address. Answer:
\begin{itemize}
  \item Do they? If your hand figure differs, which \code{ranges} cell did you use as the parent
    base?
  \item The driver's IRQ number and the GIC number from \textbf{a)} are \emph{different numbers},
    and both are right. Explain, and say which one \file{/proc/interrupts} shows in which column.
  \item Repeat \textbf{a)} with \code{od -An -tx4} instead of \code{-tx1}. Every non-zero cell
    changes. Which is correct, what is \code{-tx4} doing, and why is the wrong answer plausible
    rather than obviously broken?
\end{itemize}

\xpart{e} \textbf{The point of the whole chapter.} \code{grep} your driver's source for the address
and for the interrupt number. Neither is there.
\begin{itemize}
  \item Where is the address, then? Name the file, the node and the property.
  \item Give QEMU a \textbf{second} \code{-device qa-dev} on the command line and boot again. Two
    devices appear, and your driver probes both:
\begin{console}
qa_platform c000000.qa-dev: probed, irq 17
qa_platform c001000.qa-dev: probed, irq 18
0c000000-0c000fff : c000000.qa-dev qa-dev@0
0c001000-0c001fff : c001000.qa-dev qa-dev@1000
 17:  426  0  GIC-0 144 Level  c000000.qa-dev
 18:  420  0  GIC-0 145 Level  c001000.qa-dev
\end{console}

    Read the two devices' \code{interrupts} attributes. They differ. What would have happened if
    your private structure had been a file-scope \code{static} rather than \code{devm_kzalloc}ed per
    device? Would your Chapter~\ref{c6:ch} driver, with its \code{base=} parameter, have coped with
    this at all?
  \item Notice the second device's address. You have seen \code{0xC001000} before, in
    Exercise~\ref{c6:ex:9}, where reading it took the machine down. What is different now, and
    what does that say about how much a driver can infer from an address alone?
  \item State what would have to change in your driver to move this device to a different SoC.
\end{itemize}

\xpart{f} A colleague proposes putting the address back in the driver as a fallback, ``in case the
device tree is wrong''. Give the argument against, in terms of what a device tree is and who owns
it.
